\section{Forelesning 1: \emph{Introduksjon -- Hva er strategi?}}
\label{sec:lec01}

\subsection*{Oversikt}
Den første forelesningen etablerer det begrepsmessige fundamentet for hele kurset:
hva strategi er, hva strategi \emph{ikke} er, og hvilke antagelser enhver strategi
hviler på. Forelesningen introduserer også kursets gjennomgående linse -- at
IT-strategi handler om bevisste valg om \emph{hvordan} digital teknologi skal brukes
til å posisjonere virksomheten, ikke bare om å digitalisere eksisterende drift.

\subsection{Hovedteori}

\subsubsection*{Strategi som posisjonering (Porter 1996)}
Michael Porter sin klassiske artikkel \emph{What Is Strategy?} er det teoretiske
ankeret. Hovedpoenget er at strategi handler om å \keypoint{utføre andre aktiviteter
enn konkurrentene, eller utføre lignende aktiviteter på en annen måte}. Strategi
betyr å være \emph{annerledes}, ikke bare \emph{bedre}.

\begin{definition}{Strategisk posisjonering}
Et bevisst valg av en unik kombinasjon av aktiviteter som skaper en
verdiblanding ingen konkurrenter kan matche uten å gi opp sin egen posisjon.
\end{definition}

Porter skiller skarpt mellom \term{operasjonell effektivitet} og \term{strategi}:
\begin{itemize}
  \item \textbf{Operasjonell effektivitet} -- å gjøre de samme aktivitetene bedre
    enn konkurrentene. Konvergerer over tid (alle kopierer best practice).
  \item \textbf{Strategi} -- å velge bevisst hvilke aktiviteter man skal gjøre, og
    hvordan disse henger sammen i et unikt system.
\end{itemize}

\subsubsection*{Strategi som valg og avveininger}
Et kjernepoeng: strategi krever \keypoint{trade-offs}. Man kan ikke samtidig være
den billigste og den mest premium aktøren. Et bevisst \emph{fravalg} er like
viktig som et \emph{tilvalg}.

\subsubsection*{Antagelser bak en strategi}
Enhver strategi hviler på antagelser om:
\begin{itemize}
  \item \emph{Omgivelsene} -- konkurrenter, kunder, teknologi, regulering.
  \item \emph{Virksomhetens kapabiliteter} -- hva man kan og ikke kan.
  \item \emph{Hva verdi er} -- for hvem, og hvordan den fanges.
\end{itemize}
Strategisk refleksjon innebærer å gjøre disse antagelsene eksplisitte og teste dem.

\subsection{Sentrale fagbegreper}
\begin{itemize}
  \item \term{Strategy} -- bevisst valg av posisjon og aktivitetssystem.
  \item \term{Operational effectiveness} -- å gjøre eksisterende aktiviteter bedre.
  \item \term{Strategic positioning} -- å være annerledes, ikke bare bedre.
  \item \term{Trade-offs} -- bevisste fravalg som beskytter posisjonen.
  \item \term{Fit} -- hvordan aktiviteter henger sammen og forsterker hverandre.
\end{itemize}

\subsection{Modeller og rammeverk}

\figureplaceholder{lec01-strategy-vs-operational-effectiveness}{Porters skille
mellom strategisk posisjonering og operasjonell effektivitet.}

Et rammeverk for å spørre om en aktivitet er strategi eller drift:
\begin{enumerate}
  \item Er aktiviteten en del av et unikt system?
  \item Krever den eksplisitte fravalg?
  \item Forsterker den andre aktiviteter i virksomheten?
\end{enumerate}

\subsection{Eksempler og caser}
\begin{itemize}
  \item \textbf{Southwest Airlines} -- Porters klassiske eksempel: bevisst valg
    om kun korte ruter, ingen bagasje-overføring, ingen måltider. Hele
    aktivitetssystemet er konsistent med posisjonen \emph{lavpris, høyfrekvent}.
  \item \textbf{IT-strategi-perspektiv} -- en digital investering må vurderes
    mot om den forsterker eller undergraver den eksisterende posisjonen.
    Coop App er et eksempel: den må styrke butikk-først-posisjonen, ikke
    konkurrere med den.
\end{itemize}

\subsection{Eksamensrelevans}
\begin{itemize}
  \item Vær klar til å definere strategi presist (Porter 1996) og skille fra
    operasjonell effektivitet.
  \item Forklar hvorfor trade-offs er nødvendige.
  \item Bruk strategi-som-posisjonering når du argumenterer for at en
    IT-investering er eller ikke er strategisk (jf.\ \autoref{sec:lec04} om
    digital transformasjon vs.\ IT-enabled transformation).
  \item Kunne diskutere antagelser bak en strategi -- dette er en
    metarefleksjon som eksaminator setter pris på.
\end{itemize}
